% Meshtastic channel access, 2 of 4: the interrupts and notifications that
% drive the MAC loop.
% Build from this folder: pdflatex channel-access-irq.tex  (palette.tex).
%
% One of four figures on Meshtastic channel access, all from the same reading
% of meshtastic/firmware tag v2.7.26.54e0d8d (the firmware_target of the
% deployments): mesh/RadioLibInterface.cpp (isrLevel0Common, onNotify,
% setTransmitDelay, startSend, completeSending, handleReceiveInterrupt,
% pollMissedIrqs), mesh/SX126xInterface.cpp (startReceive), airtime.{h,cpp}.
%   channel-access-trip   — the MAC loop alone, nothing else on the page
%   channel-access-loop   — the same loop, plus what each sense actually reads
%   channel-access-irq    — the interrupts and notifications that drive the loop
%   channel-access-model  — the arithmetic behind every number
%   channel-access-window — the contention windows on a time axis
\documentclass[tikz,border=6pt]{standalone}
\usepackage{amsmath}
% Shared look for the diagrams in this folder: font, palette, TikZ styles.
% % Shared look for the diagrams in this folder: font, palette, TikZ styles.
% % Shared look for the diagrams in this folder: font, palette, TikZ styles.
% \input{palette} in the preamble of a standalone tikz document.
%
% Font: Times New Roman with lualatex/xelatex (TeX Gyre Termes, a Times clone,
% where it is not installed); Times via mathptmx with pdflatex.
%
% Palette:
%   Meshtastic  #2C2D3C dark, #67EA94 green; for lines and text on white,
%               Green 600 #3FB86D and Green 700 #2D8F52
%               (github.com/meshtastic/design).
%   LoRa        #00AEEF "LoRa blue" (Semtech LoRa logo usage guidelines).
%   Connector   amber, the one accent outside both palettes.
\usepackage{iftex}
\ifPDFTeX
  \usepackage[T1]{fontenc}
  \usepackage{mathptmx}
\else
  \usepackage{fontspec}
  \IfFontExistsTF{Times New Roman}
    {\setmainfont{Times New Roman}}
    {\setmainfont{TeX Gyre Termes}}
\fi
\usetikzlibrary{positioning,arrows.meta,fit,backgrounds,calc,
                decorations.pathmorphing,shapes.geometric}

\definecolor{mtdark}{HTML}{2C2D3C}
\definecolor{mtgreen}{HTML}{67EA94}
\definecolor{mtgreen6}{HTML}{3FB86D}
\definecolor{mtgreen7}{HTML}{2D8F52}
\definecolor{lorablue}{HTML}{00AEEF}
\definecolor{amber}{HTML}{E69500}

% the second line of a radio's label, in Meshtastic green
\newcommand{\radiosub}[1]{{\scriptsize\color{mtgreen}#1}}

\tikzset{
  >={Stealth[length=2.2mm]},
  box/.style={draw, rounded corners=2pt, align=center, inner sep=4pt, minimum height=11mm},
  % a Meshtastic radio: dark, green-edged, white title
  radio/.style={box, very thick, fill=mtdark, draw=mtgreen, text=white, minimum width=28mm},
  % what owns a box: this platform, the PBX project, the gateway
  platform/.style={box, fill=mtgreen!16, draw=mtgreen7},
  pbx/.style={box, fill=amber!12, draw=amber!85!black},
  gateway/.style={box, fill=mtdark!7, draw=mtdark!70},
  other/.style={box, fill=white, draw=mtdark!45},
  member/.style={->, very thick, draw=mtgreen6},
  conn/.style={<->, very thick, draw=amber},
  dep/.style={->, thick, dashed, draw=mtdark!70},
  air/.style={draw=lorablue, thick, decorate,
              decoration={snake, amplitude=1.3pt, segment length=6pt}},
  note/.style={font=\scriptsize, align=center},
}
 in the preamble of a standalone tikz document.
%
% Font: Times New Roman with lualatex/xelatex (TeX Gyre Termes, a Times clone,
% where it is not installed); Times via mathptmx with pdflatex.
%
% Palette:
%   Meshtastic  #2C2D3C dark, #67EA94 green; for lines and text on white,
%               Green 600 #3FB86D and Green 700 #2D8F52
%               (github.com/meshtastic/design).
%   LoRa        #00AEEF "LoRa blue" (Semtech LoRa logo usage guidelines).
%   Connector   amber, the one accent outside both palettes.
\usepackage{iftex}
\ifPDFTeX
  \usepackage[T1]{fontenc}
  \usepackage{mathptmx}
\else
  \usepackage{fontspec}
  \IfFontExistsTF{Times New Roman}
    {\setmainfont{Times New Roman}}
    {\setmainfont{TeX Gyre Termes}}
\fi
\usetikzlibrary{positioning,arrows.meta,fit,backgrounds,calc,
                decorations.pathmorphing,shapes.geometric}

\definecolor{mtdark}{HTML}{2C2D3C}
\definecolor{mtgreen}{HTML}{67EA94}
\definecolor{mtgreen6}{HTML}{3FB86D}
\definecolor{mtgreen7}{HTML}{2D8F52}
\definecolor{lorablue}{HTML}{00AEEF}
\definecolor{amber}{HTML}{E69500}

% the second line of a radio's label, in Meshtastic green
\newcommand{\radiosub}[1]{{\scriptsize\color{mtgreen}#1}}

\tikzset{
  >={Stealth[length=2.2mm]},
  box/.style={draw, rounded corners=2pt, align=center, inner sep=4pt, minimum height=11mm},
  % a Meshtastic radio: dark, green-edged, white title
  radio/.style={box, very thick, fill=mtdark, draw=mtgreen, text=white, minimum width=28mm},
  % what owns a box: this platform, the PBX project, the gateway
  platform/.style={box, fill=mtgreen!16, draw=mtgreen7},
  pbx/.style={box, fill=amber!12, draw=amber!85!black},
  gateway/.style={box, fill=mtdark!7, draw=mtdark!70},
  other/.style={box, fill=white, draw=mtdark!45},
  member/.style={->, very thick, draw=mtgreen6},
  conn/.style={<->, very thick, draw=amber},
  dep/.style={->, thick, dashed, draw=mtdark!70},
  air/.style={draw=lorablue, thick, decorate,
              decoration={snake, amplitude=1.3pt, segment length=6pt}},
  note/.style={font=\scriptsize, align=center},
}
 in the preamble of a standalone tikz document.
%
% Font: Times New Roman with lualatex/xelatex (TeX Gyre Termes, a Times clone,
% where it is not installed); Times via mathptmx with pdflatex.
%
% Palette:
%   Meshtastic  #2C2D3C dark, #67EA94 green; for lines and text on white,
%               Green 600 #3FB86D and Green 700 #2D8F52
%               (github.com/meshtastic/design).
%   LoRa        #00AEEF "LoRa blue" (Semtech LoRa logo usage guidelines).
%   Connector   amber, the one accent outside both palettes.
\usepackage{iftex}
\ifPDFTeX
  \usepackage[T1]{fontenc}
  \usepackage{mathptmx}
\else
  \usepackage{fontspec}
  \IfFontExistsTF{Times New Roman}
    {\setmainfont{Times New Roman}}
    {\setmainfont{TeX Gyre Termes}}
\fi
\usetikzlibrary{positioning,arrows.meta,fit,backgrounds,calc,
                decorations.pathmorphing,shapes.geometric}

\definecolor{mtdark}{HTML}{2C2D3C}
\definecolor{mtgreen}{HTML}{67EA94}
\definecolor{mtgreen6}{HTML}{3FB86D}
\definecolor{mtgreen7}{HTML}{2D8F52}
\definecolor{lorablue}{HTML}{00AEEF}
\definecolor{amber}{HTML}{E69500}

% the second line of a radio's label, in Meshtastic green
\newcommand{\radiosub}[1]{{\scriptsize\color{mtgreen}#1}}

\tikzset{
  >={Stealth[length=2.2mm]},
  box/.style={draw, rounded corners=2pt, align=center, inner sep=4pt, minimum height=11mm},
  % a Meshtastic radio: dark, green-edged, white title
  radio/.style={box, very thick, fill=mtdark, draw=mtgreen, text=white, minimum width=28mm},
  % what owns a box: this platform, the PBX project, the gateway
  platform/.style={box, fill=mtgreen!16, draw=mtgreen7},
  pbx/.style={box, fill=amber!12, draw=amber!85!black},
  gateway/.style={box, fill=mtdark!7, draw=mtdark!70},
  other/.style={box, fill=white, draw=mtdark!45},
  member/.style={->, very thick, draw=mtgreen6},
  conn/.style={<->, very thick, draw=amber},
  dep/.style={->, thick, dashed, draw=mtdark!70},
  air/.style={draw=lorablue, thick, decorate,
              decoration={snake, amplitude=1.3pt, segment length=6pt}},
  note/.style={font=\scriptsize, align=center},
}


\tikzset{
  mac/.style={box, fill=mtgreen!16, draw=mtgreen7, align=center, font=\small},
  case/.style={box, fill=white, draw=mtdark!55, align=left, text width=46mm, font=\scriptsize},
  low/.style={box, fill=lorablue!7, draw=lorablue!70, align=center, font=\scriptsize},
  hw/.style={->, thick, draw=lorablue!80!black},
  up/.style={->, thick, draw=mtdark!75},
  elab/.style={font=\scriptsize\itshape, inner sep=1.5pt, fill=white, text=mtdark!75},
  panel/.style={font=\bfseries\small, text=mtdark},
}
\newcommand{\head}[1]{\textbf{#1}}

\begin{document}
\begin{tikzpicture}[font=\small]

\node[panel, anchor=west] at (-8.2,1.9) {What drives the loop: the radio's interrupts};

% ── the MAC loop, seen from below ────────────────────────────────────────────
\node[mac, minimum width=150mm] (mac) at (0.6,0.8)
  {\head{the MAC loop} --- figure \texttt{channel-access-loop}\\[-1pt]
   {\scriptsize draw $\delta$ \textperiodcentered\ wait \textperiodcentered\ busy?
    \textperiodcentered\ CAD \textperiodcentered\ transmit}};

% ── the three things onNotify() can be woken with ────────────────────────────
\node[case] (onrx) at (-4.6,-2.6)
  {\head{\texttt{ISR\_RX}}\\[2pt]
   \texttt{handleReceiveInterrupt()}: the air time goes into $U_{ch}$, the frame is
   decoded and handed to the router (figure \texttt{rebroadcast}).};
\node[case] (ontx) at (0.6,-2.6)
  {\head{\texttt{ISR\_TX}}\\[2pt]
   \texttt{handleTransmitInterrupt()} $\to$ \texttt{completeSending()}: the ToA just
   transmitted goes into $U_{ch}$, the packet is freed, \texttt{txGood++}.};
\node[case] (ontim) at (5.8,-2.6)
  {\head{\texttt{TRANSMIT\_DELAY\_COMPLETED}}\\[2pt]
   The backoff elapsed: run the queue --- busy?, CAD, and send or re-draw.};

\begin{scope}[on background layer]
  \node[draw=mtdark!45, dashed, rounded corners=6pt, fill=mtdark!3,
        inner xsep=7pt, inner ysep=13pt, fit=(onrx)(ontx)(ontim)] (grp) {};
\end{scope}
\node[font=\scriptsize\bfseries, anchor=north west, text=mtdark!85] at ($(grp.north west)+(3pt,-2pt)$)
  {\texttt{onNotify(cause)} --- on the radio thread, one cause per wake-up};

% ── below: where each cause is born ──────────────────────────────────────────
\node[low, text width=66mm] (isr) at (-2.0,-6.4)
  {\head{\texttt{isrLevel0Common(cause)}} --- in interrupt context\\[2pt]
   \texttt{disableInterrupt()}, then \texttt{notifyFromISR(cause)}};

\node[low, text width=46mm] (tmr) at (5.8,-6.4)
  {\head{the backoff timer}\\[2pt]
   \texttt{notifyLater($\delta$, TRANSMIT\_DELAY\_COMPLETED)}};

\node[radio, minimum width=66mm] (rad) at (-2.0,-8.8)
  {\textbf{SX1262 \textperiodcentered\ DIO1}\\[-1pt]
   \radiosub{\texttt{RxDone} \textperiodcentered\ \texttt{TxDone} --- the only two
   events on the wire}};

\draw[hw] (rad.north) -- node[elab, right] {a rising edge on DIO1} (isr.south);
\draw[hw] (isr.north) -| node[elab, above, pos=0.25] {\texttt{ISR\_RX}} (onrx.south);
\draw[hw] (isr.north) -| node[elab, above, pos=0.25] {\texttt{ISR\_TX}} (ontx.south);
\draw[hw] (tmr.north) -- (ontim.south);

\draw[up] (onrx.north) -- node[elab, right, align=left] {\texttt{startReceive()}\\ \texttt{setTransmitDelay()}}
          (onrx.north |- mac.south);
\draw[up] (ontx.north) -- node[elab, right, align=left] {\texttt{startReceive()}\\ \texttt{setTransmitDelay()}}
          (ontx.north |- mac.south);
\draw[up] (ontim.north) -- (ontim.north |- mac.south);

\draw[hw] (mac.east) -- ++(0.7,0) |- node[elab, right, pos=0.25] {armed on every wait} (tmr.east);
\draw[hw] (mac.east) -- ++(2.3,0) |- (rad.east);
\node[elab, rotate=90, text=mtdark!75] at (9.8,-4.6)
  {\texttt{startSend()} / \texttt{startReceive()}: arm \texttt{isrTxLevel0} / \texttt{isrRxLevel0}};

% ── two notes ────────────────────────────────────────────────────────────────
\node[note, text width=76mm, align=left, anchor=north west, text=mtdark!80] at (-8.2,-10.0)
  {\textbf{This is where $U_{ch}$ is fed.} Every reception is logged, decodable or
   not: \texttt{RX\_ALL\_LOG} for a frame rejected as noise, too short, bad CRC or
   another channel, \texttt{RX\_LOG} for a valid mesh packet, \texttt{TX\_LOG} in
   \texttt{completeSending()}. Channel utilisation sums all three, so a node's own
   view of the channel includes traffic it could never decode.};

\node[note, text width=76mm, align=left, anchor=north west, text=mtdark!80] at (0.6,-10.0)
  {\textbf{Edge-triggered, and a single slot.} An edge can be missed, so
   \texttt{pollMissedIrqs()} re-reads the \texttt{RX\_DONE} flag about once a second
   and re-notifies \texttt{ISR\_RX} if it finds one. The thread has one
   notification slot and \texttt{notifyLater()} does not overwrite a pending
   notification. Between packets the chip is not listening continuously: it sniffs
   in Rx duty-cycle mode, which is why the preamble is 16 symbols long.};

\end{tikzpicture}
\end{document}
